% ============================================================
%  EXAMEN TEMA 1: OPERACIONES EN Q
%  Curso de Introducción — 2do Año
%  Duración: 90 minutos
%  Temas: Fracciones, Suma/Resta, Multiplicación/División,
%         Operaciones combinadas, Potenciación y Ecuaciones
%  Autor: Ing. Gabriel Astudillo
%  Nota: enunciado limpio (sin pistas ni desarrollo). Las
%  soluciones están en la "Clave de Respuestas" al final.
% ============================================================

\documentclass[12pt, letterpaper]{article}

% ── Paquetes ─────────────────────────────────────────────────

\usepackage{mathwizards-guia}
\mwguia{Examen — Tema 1: Operaciones en $\mathbb{Q}$}{Ing. Gabriel Astudillo $\cdot$ Curso 2do Año}

% ── Comandos propios de este examen ───────────────────
\newcommand{\Q}{\mathbb{Q}}

\begin{document}
% ============================================================

% ── Portada / Título ─────────────────────────────────────────
\mwportada{Examen del Tema 1}{Operaciones en $\mathbb{Q}$}{Fracciones, Operaciones Combinadas, Potenciación y Ecuaciones}

\vspace{4pt}
\begin{cuadroinfo}
  \small
  \textbf{Nombre:} \rule{0.55\linewidth}{0.4pt} \quad \textbf{Fecha:} \rule{0.15\linewidth}{0.4pt} \\[6pt]
  \textbf{Tiempo disponible:} 90 minutos \quad$\bullet$\quad \textbf{Puntaje total:} 100 puntos \\[4pt]
  \textbf{Instrucciones:}
  \begin{enumerate}[leftmargin=*]
    \item Lee cada ejercicio con calma antes de resolverlo.
    \item Muestra \emph{todos} tus pasos: se evalúa cómo llegas a la respuesta, no solo la respuesta.
    \item Simplifica \emph{siempre} la fracción final hasta obtener la fracción irreductible.
    \item Presta atención a la jerarquía de las operaciones y al signo de las potencias.
    \item No se permite el uso de calculadora.
    \item Escribe con letra clara y revisa tu examen antes de entregarlo.
  \end{enumerate}
\end{cuadroinfo}

\vspace{8pt}

% ============================================================
\section{Fracciones: Definición y Tipos (12 puntos)}
% ============================================================

\begin{enumerate}[leftmargin=*]
  \item \textbf{(4 pts)} Clasifica cada fracción en propia o impropia y escribe las impropias como número mixto:
        \[
          a)\ \frac{47}{6} \qquad b)\ \frac{5}{12} \qquad c)\ \frac{100}{9} \qquad d)\ \frac{23}{23}
        \]

  \item \textbf{(4 pts)} Simplifica hasta la fracción irreductible:
        \[
          a)\ \frac{252}{378} \qquad b)\ \frac{420}{588} \qquad c)\ \frac{693}{924}
        \]

  \item \textbf{(4 pts)} Ordena de menor a mayor las fracciones:
        \[
          \frac{5}{6}, \quad \frac{7}{9}, \quad \frac{11}{15}, \quad \frac{3}{4}
        \]
\end{enumerate}

\newpage

% ============================================================
\section{Suma y Resta de Fracciones (16 puntos)}
% ============================================================

\begin{enumerate}[leftmargin=*]
  \item \textbf{(4 pts)} Calcula y simplifica:
        \[
          a)\ \frac{5}{6} + \frac{7}{8} - \frac{3}{4} \qquad b)\ \frac{7}{12} - \frac{5}{18} + \frac{2}{9}
        \]

  \item \textbf{(4 pts)} Calcula (números mixtos):
        \[
          a)\ 3\frac{3}{4} - 2\frac{5}{6} \qquad b)\ 5\frac{1}{2} + 3\frac{3}{8} - 1\frac{3}{4}
        \]

  \item \textbf{(4 pts)} Un depósito contiene $\dfrac{7}{8}$ de su capacidad. Se extrae $\dfrac{1}{3}$ de lo que contiene y luego se añade $\dfrac{5}{12}$ de la capacidad total. ¿Qué fracción de la capacidad queda en el depósito?

  \item \textbf{(4 pts)} Entre tres hermanos se reparte una herencia: el mayor recibe $\dfrac{2}{5}$, el mediano $\dfrac{1}{3}$ y el menor el resto. Si el menor recibió $\$12\,000$, ¿cuál era el monto total de la herencia?
\end{enumerate}

\newpage

% ============================================================
\section{Multiplicación y División de Fracciones (16 puntos)}
% ============================================================

\begin{enumerate}[leftmargin=*]
  \item \textbf{(4 pts)} Calcula y simplifica:
        \[
          a)\ \frac{14}{15} \cdot \frac{25}{21} \qquad b)\ \frac{9}{20} \div \frac{27}{40}
        \]

  \item \textbf{(4 pts)} Calcula (números mixtos):
        \[
          a)\ 4\frac{1}{2} \div 2\frac{1}{4} \qquad b)\ 3\frac{1}{3} \cdot 1\frac{4}{5}
        \]

  \item \textbf{(4 pts)} Calcula:
        \[
          \left(\frac{3}{5} \cdot \frac{10}{9}\right) \div \left(\frac{2}{3} \div \frac{4}{9}\right)
        \]

  \item \textbf{(4 pts)} Un terreno de $\dfrac{5}{6}$ de hectárea se divide en parcelas iguales de $\dfrac{1}{12}$ de hectárea. ¿Cuántas parcelas se obtienen?
\end{enumerate}

\newpage

% ============================================================
\section{Operaciones Combinadas (20 puntos)}
% ============================================================

\begin{enumerate}[leftmargin=*]
  \item \textbf{(5 pts)} Calcula, respetando la jerarquía:
        \[
          \left[\left(\frac{2}{3} + \frac{1}{6}\right) \cdot \frac{9}{5} - \frac{1}{2}\right] \div \frac{7}{10}
        \]

  \item \textbf{(5 pts)} Calcula:
        \[
          \left(\frac{3}{4} - \frac{1}{2} \cdot \frac{5}{6}\right) \div \left(\frac{2}{3} + \frac{1}{6}\right)
        \]

  \item \textbf{(5 pts)} Calcula:
        \[
          \frac{\frac{2}{3} \cdot \frac{9}{4} - \frac{1}{6}}{\frac{5}{8} \div \frac{5}{16} + \frac{1}{2}}
        \]

  \item \textbf{(5 pts)} Una persona tiene $\dfrac{5}{6}$ de litro de jugo. Bebe $\dfrac{2}{5}$ de lo que tiene y luego regala $\dfrac{1}{3}$ de lo que queda. ¿Cuánto jugo le queda al final?
\end{enumerate}

\newpage

% ============================================================
\section{Potenciación en $\mathbb{Q}$ (20 puntos)}
% ============================================================

\begin{enumerate}[leftmargin=*]
  \item \textbf{(4 pts)} Calcula, distinguiendo el papel de los paréntesis:
        \[
          a)\ \left(-\frac{2}{3}\right)^4 \qquad b)\ -\left(\frac{2}{3}\right)^4 \qquad c)\ \left(-\frac{3}{2}\right)^3 \qquad d)\ \left(\frac{5}{4}\right)^0
        \]

  \item \textbf{(4 pts)} Calcula aplicando exponentes negativos:
        \[
          a)\ \left(\frac{5}{4}\right)^{-2} \qquad b)\ \left(-\frac{2}{3}\right)^{-3} \qquad c)\ \left(\frac{3}{7}\right)^{-1}
        \]

  \item \textbf{(4 pts)} Aplica las propiedades de las potencias y simplifica:
        \[
          a)\ \left(\frac{3}{5}\right)^4 \cdot \left(\frac{3}{5}\right)^{-2} \qquad
          b)\ \left(\frac{2}{7}\right)^5 \div \left(\frac{2}{7}\right)^2 \qquad
          c)\ \left[\left(-\frac{2}{3}\right)^2\right]^3
        \]

  \item \textbf{(4 pts)} Simplifica:
        \[
          \left(\frac{16x^4y^3}{8x^6y^7}\right)^{-1}
        \]

  \item \textbf{(4 pts)} Simplifica:
        \[
          \left(\frac{27a^5b^2}{9a^2b^6}\right)^{-2}
        \]
\end{enumerate}

\newpage

% ============================================================
\section{Ecuaciones Lineales con Fracciones (16 puntos)}
% ============================================================

\begin{enumerate}[leftmargin=*]
  \item \textbf{(4 pts)} Resuelve: $\dfrac{x}{2} + \dfrac{x}{3} + \dfrac{x}{4} = 13$

  \item \textbf{(4 pts)} Resuelve: $\dfrac{2x - 1}{3} + \dfrac{x + 2}{4} = \dfrac{5x}{6}$

  \item \textbf{(4 pts)} Resuelve: $\dfrac{x + 1}{2} - \dfrac{x - 2}{3} = \dfrac{2x - 1}{6}$

  \item \textbf{(4 pts)} La cuarta parte de un número, aumentada en $\dfrac{1}{2}$, es igual a la sexta parte del número más 2. ¿Cuál es el número?
\end{enumerate}

\newpage

% ============================================================
\ifmwclaves
  \section{Clave de Respuestas}
  % ============================================================

  \subsection*{Sección 1: Fracciones: Definición y Tipos}

  \begin{enumerate}[leftmargin=*, label=\textbf{\arabic*.}]
    \item a) impropia, $\dfrac{47}{6} = 7\dfrac{5}{6}$ \quad b) propia \quad c) impropia, $\dfrac{100}{9} = 11\dfrac{1}{9}$ \quad d) impropia, $\dfrac{23}{23} = 1$

    \item
          a) $\dfrac{252}{378} = \dfrac{252 \div 126}{378 \div 126} = \dfrac{2}{3}$ \quad
          b) $\dfrac{420}{588} = \dfrac{420 \div 84}{588 \div 84} = \dfrac{5}{7}$ \quad
          c) $\dfrac{693}{924} = \dfrac{693 \div 231}{924 \div 231} = \dfrac{3}{4}$

    \item Llevamos todo al mcm $= 180$:
          $$\frac{5}{6} = \frac{150}{180}, \quad \frac{7}{9} = \frac{140}{180}, \quad \frac{11}{15} = \frac{132}{180}, \quad \frac{3}{4} = \frac{135}{180}.$$
          Como $132 < 135 < 140 < 150$:
          $$\frac{11}{15} < \frac{3}{4} < \frac{7}{9} < \frac{5}{6}.$$
  \end{enumerate}

  \newpage

  \subsection*{Sección 2: Suma y Resta de Fracciones}

  \begin{enumerate}[leftmargin=*, label=\textbf{\arabic*.}]
    \item
          a) mcm$(6,8,4) = 24$: $\dfrac{20}{24} + \dfrac{21}{24} - \dfrac{18}{24} = \dfrac{23}{24}$ \\[2pt]
          b) mcm$(12,18,9) = 36$: $\dfrac{21}{36} - \dfrac{10}{36} + \dfrac{8}{36} = \dfrac{19}{36}$

    \item
          a) $3\dfrac{3}{4} - 2\dfrac{5}{6} = \dfrac{15}{4} - \dfrac{17}{6} = \dfrac{45}{12} - \dfrac{34}{12} = \dfrac{11}{12}$ \\[2pt]
          b) $5\dfrac{1}{2} + 3\dfrac{3}{8} - 1\dfrac{3}{4} = \dfrac{11}{2} + \dfrac{27}{8} - \dfrac{7}{4} = \dfrac{44}{8} + \dfrac{27}{8} - \dfrac{14}{8} = \dfrac{57}{8} = 7\dfrac{1}{8}$

    \item Se extrae $\dfrac{1}{3}$ de $\dfrac{7}{8} = \dfrac{7}{24}$. Queda $\dfrac{7}{8} - \dfrac{7}{24} = \dfrac{21}{24} - \dfrac{7}{24} = \dfrac{14}{24} = \dfrac{7}{12}$.
          Se añade $\dfrac{5}{12}$: $\dfrac{7}{12} + \dfrac{5}{12} = \dfrac{12}{12} = 1$. \textbf{El depósito queda lleno.}

    \item El menor recibe $1 - \left(\dfrac{2}{5} + \dfrac{1}{3}\right) = 1 - \dfrac{6 + 5}{15} = 1 - \dfrac{11}{15} = \dfrac{4}{15}$ del total.
          Si $\dfrac{4}{15}$ del total es $\$12\,000$, el total es $12\,000 \cdot \dfrac{15}{4} = \$45\,000$.
  \end{enumerate}

  \newpage

  \subsection*{Sección 3: Multiplicación y División de Fracciones}

  \begin{enumerate}[leftmargin=*, label=\textbf{\arabic*.}]
    \item
          a) $\dfrac{14}{15} \cdot \dfrac{25}{21} = \dfrac{350}{315} = \dfrac{10}{9}$ \\[2pt]
          b) $\dfrac{9}{20} \div \dfrac{27}{40} = \dfrac{9}{20} \cdot \dfrac{40}{27} = \dfrac{360}{540} = \dfrac{2}{3}$

    \item
          a) $4\dfrac{1}{2} \div 2\dfrac{1}{4} = \dfrac{9}{2} \div \dfrac{9}{4} = \dfrac{9}{2} \cdot \dfrac{4}{9} = 2$ \\[2pt]
          b) $3\dfrac{1}{3} \cdot 1\dfrac{4}{5} = \dfrac{10}{3} \cdot \dfrac{9}{5} = \dfrac{90}{15} = 6$

    \item $\dfrac{3}{5} \cdot \dfrac{10}{9} = \dfrac{30}{45} = \dfrac{2}{3}$ \quad y \quad $\dfrac{2}{3} \div \dfrac{4}{9} = \dfrac{2}{3} \cdot \dfrac{9}{4} = \dfrac{18}{12} = \dfrac{3}{2}$.
          Entonces: $\dfrac{2}{3} \div \dfrac{3}{2} = \dfrac{2}{3} \cdot \dfrac{2}{3} = \dfrac{4}{9}$.

    \item $\dfrac{5}{6} \div \dfrac{1}{12} = \dfrac{5}{6} \cdot 12 = \dfrac{60}{6} = 10$ parcelas.
  \end{enumerate}

  \newpage

  \subsection*{Sección 4: Operaciones Combinadas}

  \begin{enumerate}[leftmargin=*, label=\textbf{\arabic*.}]
    \item
          \begin{enumerate}[label=\alph*)]
            \item Paréntesis: $\dfrac{2}{3} + \dfrac{1}{6} = \dfrac{4}{6} + \dfrac{1}{6} = \dfrac{5}{6}$.
            \item Producto: $\dfrac{5}{6} \cdot \dfrac{9}{5} = \dfrac{45}{30} = \dfrac{3}{2}$.
            \item Resta: $\dfrac{3}{2} - \dfrac{1}{2} = 1$.
            \item División: $1 \div \dfrac{7}{10} = \dfrac{10}{7}$.
          \end{enumerate}
          Resultado: $\dfrac{10}{7}$.

    \item $\dfrac{1}{2} \cdot \dfrac{5}{6} = \dfrac{5}{12}$; entonces $\dfrac{3}{4} - \dfrac{5}{12} = \dfrac{9}{12} - \dfrac{5}{12} = \dfrac{4}{12} = \dfrac{1}{3}$.
          Por otro lado, $\dfrac{2}{3} + \dfrac{1}{6} = \dfrac{5}{6}$.
          Resultado: $\dfrac{1}{3} \div \dfrac{5}{6} = \dfrac{1}{3} \cdot \dfrac{6}{5} = \dfrac{2}{5}$.

    \item Numerador: $\dfrac{2}{3} \cdot \dfrac{9}{4} - \dfrac{1}{6} = \dfrac{18}{12} - \dfrac{1}{6} = \dfrac{3}{2} - \dfrac{1}{6} = \dfrac{9}{6} - \dfrac{1}{6} = \dfrac{8}{6} = \dfrac{4}{3}$.
          Denominador: $\dfrac{5}{8} \div \dfrac{5}{16} + \dfrac{1}{2} = \dfrac{5}{8} \cdot \dfrac{16}{5} + \dfrac{1}{2} = 2 + \dfrac{1}{2} = \dfrac{5}{2}$.
          Resultado: $\dfrac{4}{3} \div \dfrac{5}{2} = \dfrac{4}{3} \cdot \dfrac{2}{5} = \dfrac{8}{15}$.

    \item Bebe $\dfrac{2}{5}$ de $\dfrac{5}{6} = \dfrac{10}{30} = \dfrac{1}{3}$ L. Quedan $\dfrac{5}{6} - \dfrac{1}{3} = \dfrac{5}{6} - \dfrac{2}{6} = \dfrac{3}{6} = \dfrac{1}{2}$ L.
          Regala $\dfrac{1}{3}$ de $\dfrac{1}{2} = \dfrac{1}{6}$ L. Le quedan $\dfrac{1}{2} - \dfrac{1}{6} = \dfrac{3}{6} - \dfrac{1}{6} = \dfrac{2}{6} = \dfrac{1}{3}$ L.
  \end{enumerate}

  \newpage

  \subsection*{Sección 5: Potenciación en $\mathbb{Q}$}

  \begin{enumerate}[leftmargin=*, label=\textbf{\arabic*.}]
    \item a) $\dfrac{16}{81}$ \quad b) $-\dfrac{16}{81}$ \quad c) $-\dfrac{27}{8}$ \quad d) $1$

    \item a) $\left(\dfrac{5}{4}\right)^{-2} = \left(\dfrac{4}{5}\right)^2 = \dfrac{16}{25}$ \quad
          b) $\left(-\dfrac{2}{3}\right)^{-3} = \left(-\dfrac{3}{2}\right)^3 = -\dfrac{27}{8}$ \quad
          c) $\left(\dfrac{3}{7}\right)^{-1} = \dfrac{7}{3}$

    \item a) $\left(\dfrac{3}{5}\right)^{4-2} = \left(\dfrac{3}{5}\right)^2 = \dfrac{9}{25}$ \quad
          b) $\left(\dfrac{2}{7}\right)^{5-2} = \left(\dfrac{2}{7}\right)^3 = \dfrac{8}{343}$ \quad
          c) $\left(-\dfrac{2}{3}\right)^{2 \cdot 3} = \left(-\dfrac{2}{3}\right)^6 = \dfrac{64}{729}$

    \item $\dfrac{16x^4y^3}{8x^6y^7} = 2x^{-2}y^{-4}$. Elevando a $-1$: $\left(2x^{-2}y^{-4}\right)^{-1} = 2^{-1}x^{2}y^{4} = \dfrac{x^2y^4}{2}$.

    \item $\dfrac{27a^5b^2}{9a^2b^6} = 3a^{3}b^{-4}$. Elevando a $-2$: $3^{-2}a^{-6}b^{8} = \dfrac{b^8}{9a^6}$.
  \end{enumerate}

  \newpage

  \subsection*{Sección 6: Ecuaciones Lineales con Fracciones}

  \begin{enumerate}[leftmargin=*, label=\textbf{\arabic*.}]
    \item mcm$(2,3,4) = 12$: $6x + 4x + 3x = 156 \Rightarrow 13x = 156 \Rightarrow x = 12$.
          Comprobación: $6 + 4 + 3 = 13$ $\checkmark$

    \item mcm$(3,4,6) = 12$: $4(2x - 1) + 3(x + 2) = 10x \Rightarrow 8x - 4 + 3x + 6 = 10x \Rightarrow 11x + 2 = 10x \Rightarrow x = -2$.
          Comprobación: $\dfrac{-5}{3} + \dfrac{0}{4} = -\dfrac{5}{3}$ y $\dfrac{5(-2)}{6} = -\dfrac{10}{6} = -\dfrac{5}{3}$ $\checkmark$

    \item mcm$(2,3,6) = 6$: $3(x + 1) - 2(x - 2) = 2x - 1 \Rightarrow 3x + 3 - 2x + 4 = 2x - 1 \Rightarrow x + 7 = 2x - 1 \Rightarrow x = 8$.
          Comprobación: $\dfrac{9}{2} - \dfrac{6}{3} = \dfrac{9}{2} - 2 = \dfrac{5}{2}$ y $\dfrac{15}{6} = \dfrac{5}{2}$ $\checkmark$

    \item Sea $x$ el número: $\dfrac{x}{4} + \dfrac{1}{2} = \dfrac{x}{6} + 2$. mcm$(4,2,6) = 12$:
          $$3x + 6 = 2x + 24 \Rightarrow x = 18.$$
          El número es $18$.
  \end{enumerate}

\fi
\end{document}
